% !TeX root = ../main.tex

\begin{digest}[en]
\sloppy\emergencystretch=5em\tolerance=10000
Hydrogen energy is widely regarded as a core carrier for deep decarbonization in industrial and transportation sectors. Green hydrogen produced through water electrolysis powered by renewable energy sources, such as wind and solar power, exhibits near-zero lifecycle carbon emissions and therefore plays a pivotal role in the global transition toward carbon neutrality. Among various electrolysis technologies, Alkaline Water Electrolysis (AWE) remains the most mature and cost-effective approach for large-scale hydrogen production, owing to its long operational history, non-precious-metal catalysts, and compatibility with megawatt-level installations. As the installed capacities of wind and photovoltaic generation continue to expand rapidly, the intermittency and volatility inherent in these renewable sources pose increasingly stringent coordination challenges for multi-stack parallel AWE systems. Conventional control strategies struggle to simultaneously achieve high tracking accuracy, smooth control actions, and strict satisfaction of safety constraints under rapidly fluctuating power profiles.

Traditional Nonlinear Model Predictive Control (NMPC) has long been considered the benchmark for advanced process control due to its ability to explicitly optimize multi-objective cost functions subject to physical constraints. In the context of AWE systems, NMPC can naturally encode power tracking, temperature regulation, and actuator smoothing into a unified finite-horizon optimal control problem, while respecting hard constraints on stack temperature, cell voltage, hydrogen-to-oxygen impurity (HTO), and individual stack power. However, NMPC relies on iterative numerical optimization at every control cycle. For the four-stack AWE system considered in this thesis, a single NMPC solve takes approximately 17.6~seconds, which is clearly unacceptable for minute-level or even second-level control cycles demanded by rapidly varying renewable power. Furthermore, NMPC suffers from sensitivity to initial guesses, local optima, and complex parameter tuning, making its online deployment impractical for industrial electrolysis plants.

On the other end of the spectrum, purely data-driven methods such as deep neural networks offer extremely fast inference but lack rigorous guarantees for safety-critical constraints. A trained neural network controller essentially performs a feedforward mapping from states and references to control actions. While this mapping can approximate the NMPC policy with high fidelity under nominal conditions, it provides no certificate that the resulting actions will keep the system state within the safe operating envelope. Occasional anomalous outputs, caused by distribution shift, adversarial inputs, or inherent generalization errors, can drive the system dangerously close to or even beyond safety boundaries. For AWE systems, where temperature excursions above 90$^\circ$C or HTO concentrations exceeding 2\% can lead to catastrophic equipment damage or explosive gas mixtures, such uncontrolled behavior is unacceptable.

To bridge the gap between computational efficiency and safety guarantees, this thesis proposes a safe learning-based predictive control framework specifically tailored for four-stack parallel AWE systems. The framework is built upon three interconnected pillars: high-fidelity physics-based modeling, generative policy learning via diffusion models, and rigorous safety enforcement through Control Barrier Functions (CBF).

The first pillar establishes a comprehensive simulation model that captures the essential multi-physics coupling within the AWE system. The state vector comprises 13 variables, including heat-exchanger lye outlet temperature, four stack temperatures, separator temperature, cooling water outlet temperature, hydrogen moles in the anode of each stack, hydrogen in the separator liquid phase, and hydrogen in the separator gas phase. The control input vector consists of 9 variables: four stack currents, four lye flow rates, and a single coolant flow rate. The dynamics integrate three sub-models. The electrochemical model computes cell voltage, Faradaic efficiency, and reaction heat power as functions of current and temperature, capturing the non-affine coupling between current and Faradaic efficiency that is crucial for accurate safety analysis. The thermal model describes energy balances in each stack and the separator, accounting for convective heat exchange with lye flows, radiative heat loss, and cooling water heat removal via a logarithmic-mean-temperature-difference heat exchanger. The HTO transport model tracks hydrogen migration from the anode through the diaphragm into the oxygen-rich cathode and separator, driven by Fick diffusion and Darcy convection, ultimately determining the gas purity in the separator gas phase. The continuous-time ordinary differential equations are numerically integrated using the RK45 method with a simulation step of 0.2~s, ensuring both accuracy and stability for the stiff thermal dynamics.

The second pillar addresses the computational bottleneck of NMPC through generative policy learning. NMPC is employed offline as an expert controller to generate high-quality state-action trajectories. Two complementary classes of power profiles are used to ensure comprehensive coverage of the operating envelope. The first class comprises realistic wind power curves, reflecting the continuous yet stochastic fluctuations that the AWE system would encounter in actual renewable hydrogen plants. The second class consists of multi-platform step signals superimposed with random noise, deliberately creating large tracking errors and thermal transients that wind data alone rarely produce. This dual-class approach ensures that the learned policy is exposed to both typical operational regimes and extreme perturbations, thereby enhancing robustness.

From these expert trajectories, a Temporal Convolutional Network (TCN) Diffusion policy is trained. Unlike conventional behavior cloning, which deterministically regresses from observations to a single action, diffusion models learn the conditional probability distribution of expert actions. The training process follows the Denoising Diffusion Probabilistic Model (DDPM) framework. During the forward process, Gaussian noise is progressively added to expert actions until the signal becomes approximately pure noise. During the reverse process, a noise-prediction network conditioned on system states and power plans learns to iteratively denoise corrupted actions and recover high-quality control decisions. The TCN backbone leverages dilated causal convolutions with Mish activations across four temporal resolution levels, enabling effective aggregation of historical dynamics without the recurrent-state bottlenecks of LSTMs. Once trained, the policy generates 64 candidate action sequences in parallel during each control cycle. Each candidate is evaluated through a physically informed cost function that weights power tracking, temperature deviation, current smoothness, and actuator effort. Cost-weighted selection then picks the best candidate, and only the first action is executed, following the receding-horizon principle. This sampling-based optimization achieves inference times of approximately 177~ms, reducing the online computational burden by nearly two orders of magnitude compared to NMPC, while maintaining comparable power-tracking accuracy.

The third pillar provides the critical safety guarantee that purely learned policies lack. Control Barrier Functions (CBF) offer a rigorous theoretical framework for enforcing forward invariance of safe sets. For the AWE system, two categories of safety constraints are considered: temperature limits and HTO limits. The temperature constraints are formulated as first-order CBF conditions because the control inputs (current and lye flow) directly appear in the first Lie derivative of the temperature barrier function, yielding a relative degree of one. The HTO constraint is similarly formulated as a first-order CBF condition, leveraging the affine dependence of hydrogen transport dynamics on lye flow rate and the non-affine coupling of current through Faradaic efficiency. At each control instant, the CBF safety filter solves a nonlinear constrained optimization problem that minimally adjusts the nominal control action from the diffusion policy while ensuring that all CBF conditions are satisfied. Slack variables with large quadratic penalties handle temporary infeasibility during extreme transients, and CasADi with the IPOPT solver provides automatic differentiation and efficient non-convex optimization. Single CBF solves consistently complete within 20~ms, making the safety filter fully compatible with real-time control requirements.

Extensive closed-loop experiments validate the proposed framework under two representative scenarios. In the step-response scenario, the system initially operates at a 3000~A steady state; at $t=60$~min, the reference current steps down to 1000~A, emulating a sudden renewable power drop. The CBF filter triggers when HTO rises toward the 2\% limit, increasing lye flow and moderately adjusting current to suppress hydrogen accumulation. HTO peaks at 1.99\%, safely below the threshold, while the mean current change magnitude remains below 3~A, demonstrating smooth control behavior. In the high-dynamic wind-power scenario, the TCN Diffusion policy tracks the fluctuating reference with an RMSE of 0.14~MW and a temperature MAE of 2.0~K, closely matching NMPC performance. The CBF filter achieves 100\% constraint satisfaction throughout the entire test horizon, with no temperature, voltage, power, or HTO violations.

In summary, this thesis establishes an integrated safe learning control framework that combines the computational efficiency of generative diffusion models with the rigorous safety guarantees of Control Barrier Functions. The framework addresses the real-time coordination challenge of multi-stack AWE systems under volatile renewable power and provides a scalable methodology for safety-critical control in non-affine nonlinear chemical processes. Future research directions include Sim-to-Real transfer through domain randomization and online system identification, adaptive policy updates to accommodate equipment aging,~multi\-timescale scheduling-control co-optimization, and extension of the non-affine CBF framework to fuel cells and catalytic reactors.
\end{digest}
